\section{Linux的目录结构}

\begin{frame}
    \frametitle{Linux的目录结构}
    Linux系统的目录结构遵循文件系统层次标准（Filesystem Hierarchy Standard，FHS）。这实际上是几乎所有类Unix系统遵循的规范：你甚至可以在MacOS上找到相似的目录层次结构！
\end{frame}

\begin{frame}
    \begin{Table}[Linux的根目录结构]!
        \begin{tblr}{colspec={X[1]X[2]X[1]X[2]},cell{2-Z}{1,3}={font=\ttfamily}}
            目录&用途&目录&用途\\
            /bin   & 用户执行程序 &
                                        /media & 自动挂载点 \\
            /lib   & 动态与静态库 &
                                        /mnt   & 手动挂载点 \\
            /sbin  & 系统执行程序 &
                                        /boot  & 启动引导 \\
            /home  & 普通用户家目录 &
                                        /dev   & 设备文件 \\
            /root  & 超级用户家目录&
                                        /sys   & 设备属性 \\
            /etc   & 配置文件 &
                                        /proc  & 进程信息 \\
            /var   & 可变文件 &
                                        /run   & 进程数据 \\
            /usr   & 系统资源 &
                                        /tmp   & 临时数据 \\
            /opt   & 可选软件 & 
                                        /srv   & 服务数据 \\
        \end{tblr}
    \end{Table}
\end{frame}

\begin{frame}
    \frametitle{目录\code{/bin /lib /sbin}}
    \code{/bin}目录存放了所有用户可以执行的命令
    \begin{Code}*
        \lstinputlisting[firstline=1,lastline=2]{code/fhs/bin.sh}
    \end{Code}
    \code{/bin}目录也会包含大部分通过\code{apt}安装的程序。对于Linux系统，它不会区分所谓“系统自带的基本命令”和“用户安装的应用程序”，因为它们两者本质上都是可以执行的二进制程序。
\end{frame}

\begin{frame}
    \frametitle{目录\code{/bin /lib /sbin}}
    例如，安装\code{build-essential}包，这包含了C/C++开发的基本工具
    \begin{Code}*
        \lstinputlisting[firstline=4,lastline=4]{code/shell/apt.sh}
    \end{Code}
    你现在就可以在\code{/bin}下找到\code{gcc}和\code{make}等工具了
    \begin{Code}*
        \lstinputlisting[firstline=3,lastline=4]{code/fhs/bin.sh}
    \end{Code}
\end{frame}

\begin{frame}
    \frametitle{目录\code{/bin /lib /sbin}}
    \code{/lib}目录存放了所有动态库和静态库 
    \begin{itemize}
        \item 库的本质是预先编译好的通用二进制片段，它可以和应用程序进行链接。
        \item 动态链接库：运行时链接，后缀为\code{.so}，例如C库\code{libc.so}
        \item 静态链接库：编译时链接，后缀为\code{.a}，例如C库\code{libc.a}
    \end{itemize}
    你可以在\code{/lib/x86_64-linux-gnu}下找到\code{libc.so}和\code{libc.a}
    \begin{Code}*
        \lstinputlisting[firstline=1,lastline=2]{code/fhs/lib.sh}
    \end{Code}
\end{frame}

\begin{frame}
    \frametitle{目录\code{/bin /lib /sbin}}
    通常，动态链接是较为常用的。
    
    当你在C程序中写下\code{\#include <stdio.h>}以及\code{printf}并编译时，函数\code{printf}的二进制并不会被链接到你的程序中，真正发生的事情是，当你的程序运行到\code{printf}函数时，系统会找到\code{libc.so}并运行对应的二进制，这就是所谓“运行时链接”的含义。动态链接机制使系统中海量的二进制程序不必重复包含C库的片段，而是可以共享一个C库，从而节约大量空间！

\end{frame}

\begin{frame}
    \frametitle{目录\code{/bin /lib /sbin}}
    
    然而，静态链接在某些情况下也有其意义。
    
    当你的C程序最终是在某个裸机的单片机上运行时，根本没有操作系统进行动态链接，此时就需要在编译阶段静态链接用到的函数。空间上也很划算，只需要带上C库中真正用到的部分！
\end{frame}

\begin{frame}
    \frametitle{目录\code{/bin /lib /sbin}}
    \code{/sbin}目录存放了所有\code{root}才能执行的命令
    \begin{Code}*
        \lstinputlisting{code/fhs/sbin.sh}
    \end{Code}
    这是因为，在多人使用的服务器上，普通用户显然不能随意关机和重启！ 
\end{frame}

\begin{frame}
    \frametitle{目录\code{/home /root}}
    \code{/home}和\code{root}都是用户的家目录
    \begin{itemize}
        \item \code{/home}是普通用户的家目录，每个用户在\code{/home}下有同名的家目录。
        \item 用户\code{liyuxuan}的家目录位于\code{/home/liyuxuan}。
        \item 用户的家目录始终可以用\code{\~}表示（取决于当前登录的用户）。
        \item \code{/root}是超级用户的家目录，这是\code{root}作为超级用户在架构上享有的额外特权：其家目录并非位于通常的\code{/home/root}，而是相应地高挂了一级，直接位于根目录\code{/root}下。
    \end{itemize}
\end{frame}

\begin{frame}
    \frametitle{目录\code{/etc}}
    \code{/etc}下存放配置文件，比较典型的例子是\code{hosts}文件
    \begin{Code}*
        \lstinputlisting{code/fhs/etc.sh}
    \end{Code}
    你可以在这里设置本地的DNS解析，例如
    \begin{Code}*
        \lstinputlisting{code/hosts}
    \end{Code}
\end{frame}

\begin{frame}
    \frametitle{目录\code{/var}}
    \code{/var}下存放可变文件，比较典型的例子是\code{apt}的日志
    \begin{Code}*
        \lstinputlisting{code/fhs/var.sh}
    \end{Code}
    你可以在两个文件中分别看到\code{apt}最近的更新记录和当时的命令行输出内容。
\end{frame}

\begin{frame}
    \frametitle{目录\code{/usr}}
    \code{/usr}的全称是Unix System Resource，请注意其和User无关！

    如果你足够细心，仔细观察了\code{ls -l /}的输出，你会注意到\code{/bin /lib /sbin}实际上分别是指向\code{/usr/bin /usr/lib /usr/sbin}的软链接！这是因为早期的实践中，前者存储系统所需的命令和库，后者存储用户安装的命令和库，但是，随着时代发展，现在主流的Linux发行版都选择将这两个目录用软链接的方式合并在一起，不再区分到底是系统所需还是用户安装。
\end{frame}

\begin{frame}
    \frametitle{目录\code{/opt}}
    \code{/opt}用于安装可选的第三方大型软件，例如
    \begin{itemize}
        \item Wolfram Mathematica
        \item Xilinx Vivado
        \item 腾讯会议
    \end{itemize}
    这些软件不遵守标准的目录结构，因此统一在\code{/opt}下为每个软件分配独立的目录。
\end{frame}

\begin{frame}
    \frametitle{目录\code{/media /mnt}}
    \code{/media}和\code{/mnt}分别是自动挂载点和手动挂载点，例如当插入U盘后
    \begin{Code}*
        \lstinputlisting{code/fhs/media.sh}
    \end{Code}
    假如U盘的名称是\code{udisk}，那么你就会在该位置看到出现了一个名为\code{udisk}的目录，该过程被称为挂载（Mount），这一点和Windows中插入U盘会出现新的盘符的行为是不太一样的。
    
    通常移动存储设备都可以被自动挂载至\code{/media}，极少数情况需要手动挂载至\code{/mnt}。
\end{frame}

\begin{frame}
    \frametitle{目录\code{/boot}}
    \code{/boot}目录包含了操作系统正常启动所必须的文件。事实上，当在安装系统进行分区时，你可能会注意到除了被挂载于\code{/}的主分区，还会自动产生一个被挂载于\code{/boot/efi}的启动分区。
\end{frame}

\begin{frame}
    \frametitle{目录\code{/dev}}
    \code{/dev}用于存放设备文件，例如
    \begin{Code}*
        \lstinputlisting{code/fhs/dev.sh}
    \end{Code}
    第一条命令的输出可能会看到\code{nvme0n1p1}和\code{nvme0n1p2}，这对应了在NVMe上连接的SSD的第一个分区（\code{p1}）和第二个分区（\code{p2}）。第二条命令用于查看磁盘占用，观察该命令的输出可以发现，挂载点\code{/}对应的是\code{/dev/nvme0n1p1}，挂载点\code{/boot/efi}对应的是\code{/dev/nvme0n1p2}。
\end{frame}

\begin{frame}
    \frametitle{目录\code{/dev}}
    在Linux系统中，\alert{一切皆文件}。磁盘设备本身也被视为一个文件，这并不是一个真实存在的文件，而是通过文件读写的方式包装了对磁盘设备的数据写入和读取。而挂载的作用，就是将磁盘设备的直接读写操作，转换为我们在文件系统层面所看到的对文件的创建、修改、删除。

    设备文件在\code{ls -l}下会显示为\code{b}或\code{c}类型，代表这是一个块设备或字符设备。
\end{frame}

\begin{frame}
    \frametitle{目录\code{/sys}}
    \code{/sys}用于配置设备属性，例如
    \begin{Code}*
        \lstinputlisting{code/fhs/sys.sh}
    \end{Code}
    第一条命令将关闭键盘灯，第二条命令将开启键盘灯。
    
    这里\code{tee}最终发挥的作用就是将\code{0}和\code{1}写入\code{brightness}文件（重定向\code{>}在这里会有权限问题）。\code{/sys}下的文件是设备属性的反映，通过对\code{/sys}下文件的读写可以查看和控制设备属性。
\end{frame}

\begin{frame}
    \frametitle{目录\code{/proc}}
    \code{/proc}是进程的映射，例如
    \begin{Code}*
        \lstinputlisting{code/fhs/proc.sh}
    \end{Code}
    第一条命令会看到大量数字命名的目录，例如\code{248 3533 55991}。第二条命令用于查看全部进程。通过对照，我们可以发现，进程号和\code{/proc}下的数字命名的目录是一一对应的！
\end{frame}

\begin{frame}
    \frametitle{目录\code{/run /tmp /srv}}
    \code{/run}用于存放进程的数据，\code{/tmp}用于存放任何临时文件，系统重启后自动清空。 

    \code{/srv}用于存放本机对外提供的服务的数据（适用于服务器）。
\end{frame}
